\documentclass[11pt,a4paper]{article}
\usepackage[utf8]{inputenc}
\usepackage[T1]{fontenc}
\usepackage[spanish]{babel}
\usepackage{amsmath,amssymb,mathtools}
\usepackage{geometry}
\usepackage{enumitem}
\usepackage{array}
\usepackage{lmodern}
\usepackage{microtype}
\geometry{margin=2.4cm}
\setlength{\parindent}{0pt}
\setlength{\parskip}{0.45em}

\newcommand{\R}{\mathbb{R}}
\newcommand{\C}{\mathbb{C}}
\newcommand{\K}{\mathbb{K}}
\newcommand{\Ker}{\operatorname{Ker}}
\newcommand{\Ima}{\operatorname{Im}}
\newcommand{\rg}{\operatorname{rg}}
\newcommand{\tr}{\operatorname{tr}}
\newcommand{\Id}{\operatorname{Id}}
\newcommand{\diag}{\operatorname{diag}}
\newcommand{\Span}{\operatorname{L}}

\begin{document}

\begin{center}
{\LARGE\bfseries Álgebra Lineal}\\[0.4em]
{\Large Tema 3: Endomorfismos. Diagonalización --- Ejercicios}
\end{center}

\section{Ejercicio 1}
Si una matriz $A$ de orden $4\times4$ tiene como autovalor $\lambda=8$:
\begin{enumerate}[label=\alph*)]
\item ¿Qué posibilidades hay para $\rg(A-8I)$?
\item ¿Qué posibilidades hay para $\dim S(8)$?
\end{enumerate}

\section{Ejercicio 2}
Si una matriz $M$ de orden $2\times2$ es diagonalizable y sus autovalores son $3$ y $0$, con autovectores correspondientes $(4,1)$ y $(1,1)$, obtener $M$.

\section{Ejercicio 3}
Sabiendo que los autovalores de una matriz $A$ son $1,-1,2$, hallar los autovalores de:
\[
4A,\qquad A^t,\qquad A^3,\qquad A^{-1}.
\]

\section{Ejercicio 4}
Estudiar si son diagonalizables:
\begin{itemize}
\item Una matriz $2\times2$ con autovalores $3$ y $5$ simples.
\item Una matriz $3\times3$ con autovalores $6,7,8$ simples.
\item Una matriz $3\times3$ con autovalor $-3$ doble.
\item Una matriz $3\times3$ con autovalores $9$ doble y $1$ simple.
\end{itemize}

\section{Ejercicio 5}
Sea $f:V\to V$, donde $V$ es un espacio vectorial de dimensión finita $n$ sobre un cuerpo $\K$. Demostrar que si $B$ es una base de $V$ formada por autovectores de $f$, entonces $M_B(f)$ es diagonal.

\section{Ejercicio 6}
Sea $A$ una matriz $n\times n$ y $P(\lambda)$ su polinomio característico. Demostrar:
\begin{enumerate}[label=\alph*)]
\item Si $n=2$:
\[
P(\lambda)=\lambda^2-(\tr A)\lambda+\det A.
\]
\item Si $n=3$:
\[
P(\lambda)
=
-\lambda^3
+
(\tr A)\lambda^2
-
(A_{11}+A_{22}+A_{33})\lambda
+
\det A,
\]
donde $A_{ii}$ es el adjunto del elemento $a_{ii}$.
\end{enumerate}

\section{Ejercicio 7}
Demostrar que si $A$ es una matriz de orden $n$ con:
\[
p(\lambda)=(a-\lambda)^n,
\]
entonces $A$ es diagonalizable si y sólo si:
\[
A=aI_n.
\]

\section{Ejercicio 8}
Probar que una matriz cuadrada $A$ y su traspuesta tienen los mismos autovalores. ¿Tienen también los mismos autovectores?

\section{Ejercicio 9}
Estudiar si:
\[
A=
\begin{pmatrix}
-1&-7&1\\
0&4&0\\
-1&13&-3
\end{pmatrix},
\qquad
B=
\begin{pmatrix}
1&-3&3\\
3&-5&3\\
6&-6&4
\end{pmatrix}
\]
son diagonalizables en $\R$. En caso afirmativo, calcular la forma diagonal y la matriz de paso.

\section{Ejercicio 10}
Estudiar si es diagonalizable:
\[
A=
\begin{pmatrix}
1&0&-1&0\\
2&-1&-3&0\\
1&0&-1&0\\
0&0&0&2
\end{pmatrix}.
\]

\section{Ejercicio 11}
Estudiar si es diagonalizable:
\[
A=
\begin{pmatrix}
1&4&-2\\
0&3&0\\
1&1&1
\end{pmatrix}.
\]

\section{Ejercicio 12}
Estudiar si es diagonalizable:
\[
A=
\begin{pmatrix}
0&1&0&0\\
0&0&1&0\\
0&0&0&1\\
0&0&0&0
\end{pmatrix}.
\]

\section{Ejercicio 13}
Sea el endomorfismo de $\R^3$:
\[
f(x,y,z)=(3x,-y+az,3x+bz).
\]
¿Para qué valores de $a$ y $b$ es diagonalizable?

\section{Ejercicio 14}
En $\R^4$ se considera un endomorfismo $f$ tal que:
\[
\Ker(f-I)^3
=
\{x:x_1-x_2+x_3-x_4=0\},
\]
\[
\Ker(f-I)^2
=
\{x:x_1-x_2+x_3=0,\;x_4=0\},
\]
\[
\Ker(f-I)
=
\{x:x_1+x_3=0,\;x_4=0,\;x_2=0\},
\]
\[
\Ker f
=
\{x:x_1=x_2=x_3=0\}.
\]
\begin{enumerate}[label=\alph*)]
\item Determinar una base $B'$ tal que $M_{B'}(f)$ sea la forma canónica de Jordan.
\item Calcular la matriz de $f$ en la base canónica y su expresión analítica.
\item Determinar los autovalores de $(A-2I)^{-2}$.
\end{enumerate}

\section{Ejercicio 15}
Estudiar para qué valores de $t$ es diagonalizable en $\R$:
\[
A=
\begin{pmatrix}
t+3&t^2-10\\
1&t+1
\end{pmatrix}.
\]

\section{Ejercicio 16}
Sea:
\[
A=
\begin{pmatrix}
1&1&1&1\\
1&1&-1&-1\\
1&-1&1&-1\\
1&-1&-1&1
\end{pmatrix}.
\]
\begin{enumerate}[label=\alph*)]
\item Demostrar que $A$ es diagonalizable y hallar una matriz de paso.
\item Diagonalizar $A^2$ y $A^{-1}$.
\end{enumerate}

\section{Ejercicio 17}
Sea:
\[
M=
\begin{pmatrix}
3&-1&0\\
6&-3&2\\
8&-6&5
\end{pmatrix}.
\]
\begin{enumerate}[label=\alph*)]
\item Demostrar que $M$ no es diagonalizable en $\R$.
\item Demostrar que es diagonalizable en $\C$ y hallar una matriz de paso.
\end{enumerate}

\section{Ejercicio 18}
Sea $f$ un endomorfismo de $\R^6$ cuya matriz $A$ verifica:
\[
\rg A=1.
\]
\begin{enumerate}[label=\alph*)]
\item ¿Qué se puede decir de sus autovalores?
\item ¿Existe algún caso en que $f$ sea diagonalizable?
\item Encontrar las posibles formas de Jordan.
\end{enumerate}

\section{Ejercicio 19}
Razonar si existe un endomorfismo de $\R^6$ con un único autovalor real $\lambda$ de multiplicidad algebraica $6$ tal que:
\[
\text{a)}\quad
\rg(A-\lambda I)=4,\;
\rg(A-\lambda I)^2=2,\;
\rg(A-\lambda I)^3=1,\;
\rg(A-\lambda I)^4=0,
\]
\[
\text{b)}\quad
\rg(A-\lambda I)=3,\;
\rg(A-\lambda I)^2=2,\;
\rg(A-\lambda I)^3=1,\;
\rg(A-\lambda I)^4=0.
\]
En caso afirmativo, determinar la forma de Jordan y describir la construcción de una base de Jordan.

\section{Ejercicio 20}
¿Para qué valores de $a,b,c,d,e,f$ la matriz:
\[
A=
\begin{pmatrix}
1&a&d\\
2&b&e\\
3&c&f
\end{pmatrix}
\]
tiene como autovectores:
\[
(1,0,1),\qquad(-1,1,0),\qquad(0,1,-1)?
\]
¿Qué condición deben verificar $a$ y $b$ para que el endomorfismo admita una diagonalización ortogonal?

\section{Ejercicio 21}
Una ciudad $A$ es de tránsito. Al final de cada año, $2/3$ de su población emigra a una región $B$ y $1/3$ a una región $C$. Durante ese mismo año, $1/3$ de la población de $B$ y $1/3$ de la de $C$ se trasladan a $A$. Calcular las poblaciones en régimen estacionario sabiendo que inicialmente:
\[
(A,B,C)=(60,200,300).
\]

\section{Ejercicio 22}
En $\R^3$ se considera:
\[
A=
\begin{pmatrix}
5&0&0\\
0&-1&b\\
3&0&a
\end{pmatrix}.
\]
\begin{enumerate}[label=\alph*)]
\item Determinar para qué valores de $a$ y $b$ es diagonalizable.
\item Para $a=-1$, $b=-1$, determinar una base de Jordan y la forma canónica de Jordan.
\item Para $a=-1$, $b=-1$, utilizar Cayley--Hamilton para expresar $A^9$ en función de $A$.
\end{enumerate}

\section{Ejercicio 23}
¿Existe alguna base de $\mathcal P_2(x)$ formada por autovectores del endomorfismo cuya matriz en $\{1,x,x^2\}$ es:
\[
M=
\begin{pmatrix}
1&2&3\\
\frac12&1&\frac32\\
\frac13&\frac23&1
\end{pmatrix}?
\]
En caso afirmativo, dar una base formada por autovectores.

\section{Ejercicio 24}
Sea el endomorfismo de $\mathcal P_2(x)$ cuya matriz en $\{1,x,x^2\}$ es:
\[
M=
\begin{pmatrix}
0&3&9\\
\frac13&0&3\\
\frac19&\frac13&0
\end{pmatrix}.
\]
\begin{enumerate}[label=\alph*)]
\item Probar que es diagonalizable.
\item Encontrar una base de autovectores de $\mathcal P_2(x)$.
\item Determinar $M^{-1}$ mediante Cayley--Hamilton.
\item Calcular $M^p$ para cualquier $p\in\mathbb N$.
\end{enumerate}

\section{Ejercicio 25}
Calcular $A^{-1}$ mediante Cayley--Hamilton:
\[
A=
\begin{pmatrix}
2&0&3&-5\\
1&2&-7&-5\\
0&0&-1&0\\
0&0&1&-1
\end{pmatrix}.
\]

\section{Ejercicio 26}
Sea:
\[
M=
\begin{pmatrix}
4&2\\
3&3
\end{pmatrix}.
\]
Calcular $M^n$:
\begin{enumerate}[label=\alph*)]
\item Por diagonalización.
\item Por Cayley--Hamilton.
\end{enumerate}

\section{Ejercicio 27}
Sea:
\[
A=
\begin{pmatrix}
-14&25\\
-9&16
\end{pmatrix}.
\]
Calcular:
\[
\lim_{n\to\infty}\frac1nA^n.
\]

\section{Ejercicio 28}
Obtener la forma reducida de Jordan y la base correspondiente de:
\[
M=
\begin{pmatrix}
4&0&0&-1&1\\
-1&3&0&1&0\\
1&0&3&0&0\\
1&0&0&2&1\\
0&0&0&0&3
\end{pmatrix}.
\]

\section{Ejercicio 29}
Obtener la forma reducida de Jordan y la base correspondiente de:
\[
M=
\begin{pmatrix}
-2&1&0&-1\\
0&-2&0&4\\
-4&5&2&-4\\
0&0&0&2
\end{pmatrix}.
\]

\section{Ejercicio 30}
Hallar la forma reducida de Jordan según los valores de $\alpha$:
\[
M=
\begin{pmatrix}
1&-1&1\\
1&1&0\\
1&-\alpha&1+\alpha
\end{pmatrix}.
\]

\section{Ejercicio 31}
Obtener la forma canónica de Jordan y la matriz de paso según $k$:
\[
A=
\begin{pmatrix}
1&1&0&0\\
2k&1+k&-k&k\\
2k&1+k&1-k&k\\
0&-1&0&1
\end{pmatrix}.
\]

\section{Ejercicio 32}
Sea $A$ real de orden $6$ con autovalores:
\[
1\quad\text{con multiplicidad algebraica }2
\text{ y geométrica }1,
\]
\[
2+i\quad\text{con multiplicidad algebraica }2
\text{ y geométrica }1.
\]
Obtener su forma real de Jordan.

\section{Ejercicio 33}
Obtener la forma real de Jordan y una base correspondiente para:
\[
M=
\begin{pmatrix}
1&1&1&1\\
-2&-1&0&-1\\
0&0&-1&-1\\
0&0&2&1
\end{pmatrix}.
\]

\section{Ejercicio 34}
Dadas:
\[
A=
\begin{pmatrix}
-2&0&0\\
0&2&0\\
0&1&2
\end{pmatrix},
\qquad
B=
\begin{pmatrix}
3&0&0\\
1&3&0\\
0&1&3
\end{pmatrix},
\]
calcular:
\[
e^A,\qquad e^{At},\qquad e^B,\qquad e^{Bt}.
\]

\section{Ejercicio 35}
Dada:
\[
M=\lambda I_n,
\]
calcular:
\[
e^M.
\]

\section{Ejercicio 36}
Sea:
\[
A=
\begin{pmatrix}
0&0&\cdots&0&0\\
1&0&\cdots&0&0\\
0&1&\ddots&0&0\\
\vdots&&\ddots&\ddots&\vdots\\
0&0&\cdots&1&0
\end{pmatrix}.
\]
Calcular $e^A$.

\section{Ejercicio 37}
Sea:
\[
A=
\begin{pmatrix}
a&b\\
-b&a
\end{pmatrix}.
\]
Calcular $e^A$.

\section{Ejercicio 38}
Sea $A\in M_3(\R)$ con autovalores:
\[
\frac12\quad\text{con multiplicidad algebraica }1,
\qquad
1\quad\text{con multiplicidad algebraica }2.
\]
Calcular mediante Cayley--Hamilton:
\[
2A^4-7A^3+9A^2-5A+I.
\]

\section{Ejercicio 39}
Sea $B$ una matriz $3\times3$ con autovalores $0,1,2$. Razonar si es posible calcular:
\begin{enumerate}[label=\alph*)]
\item $\rg B$.
\item $\det(B^tB)$.
\item Los autovalores de $B^tB$.
\item Los autovalores de $(B^2+I)^{-1}$.
\end{enumerate}

\end{document}
